
%%%%%%%%
\begin{frame}{Le problème étudié}


\figSlideWithSize{exec-optim-1}{12cm}

\vfill

\begin{tabular}{p{9cm}p{5cm}}
Une requête SQL est \alert{déclarative}. Elle ne dit pas \alert{comment} calculer le résultat.

\medskip

Nous avons besoin d'une \alert{forme opératoire}\,: un programme.
&
~
\end{tabular}

    
\end{frame}


%%%%%%%%
\begin{frame}{La notion de plan d'exécution}

\figSlideWithSize{exec-optim-2}{12cm}
\vfill
\begin{tabular}{p{9cm}p{5cm}}
Dans un SGBD le programme qui exécute une requête est appelé \alert{plan d'exécution}.

\medskip
Il a une forme particulière\,: c'est un \alert{arbre}, constitué \alert{d'opérateurs}.
&
~
\end{tabular}
\end{frame}


%%%%%%%%
\begin{frame}{De la requête SQL au plan d'exécution}

\figSlideWithSize{exec-optim-3}{12cm}
\vfill

\begin{tabular}{p{9cm}p{5cm}}
Deux étapes: 
\begin{redlist}
 \item (A) plan d'exécution \alert{logique} (l'algèbre) ;
  \item (B) plan d'exécution \alert{physique} (opérateurs).
\end{redlist}
\medskip
Le SGBD s'appuie  sur un \alert{catalogue} d'opérateurs.
&
~
\end{tabular}
\end{frame}

%%%%%%%%
\begin{frame}{En quoi consiste l'optimisation\,?}


\vfill
\figSlideWithSize{exec-optim-4}{12cm}
\vfill

\alert{À chaque étape, plusieurs choix}. Le système les évalue et choisit le \guill{meilleur}.

\end{frame}

%%%%%%%%%%%%%%%%%%%%%%%%%%%%%%%%%%%%%%%%%%%%%%%%%%%%%%%
\begin{frame}{À retenir}



\begin{redbullet}
  \item SQL est conçu comme un langage qui ne détermine
     pas la manière dont une requête est exécutée (\alert{même}
     dans sa forme algébrique)

   \item Les opérations à appliquer par le système sont réduites
        à celles de l'algèbre relationnelle
        
    \item Chaque opérateur peut être évalué avec des algorithmes différents
\end{redbullet}

\vfill

Au niveau du réglage de l'organisation physique, on peut améliorer
(``optimiser'') les performances: c'est le boulot de l'administrateur
\vfill

\alert{Conclusion\,: il faut écrire les requêtes pour qu'elles soient lisibles
et claires}. L'efficacité est un problème indépendant.

\end{frame}
